%!TEX root = ../thesis.tex
\chapter{绪论}

\section{工程概述}

\subsection{工程背景和意义}
近年来，随着人工智能、物联网、大数据以及边缘计算等前沿技术的迅猛发展，机器人技术正以前所未有的速度渗透进人们的日常生活与工作场景。家用、学校等相对较小且环境复杂的场所对机器人的搜索、定位和决策提出了更高要求，而单个机器人往往因硬件资源、传感器精度和算法实时性等方面的限制，难以独立完成高效的任务。因此，多机器人协同搜索系统通过多台机器人之间的紧密配合与任务分解，有效弥补了单机在覆盖范围和处理速度上的不足，为室内安防、智能巡检以及教育科研等领域提供了值得参考的解决方案。%

在这种背景下，ROS（机器人操作系统）凭借其开源、模块化和高度可扩展的特性，成为实现多机器人系统协同作业的重要平台。ROS不仅提供了丰富的通信机制，如Topic、Service和Action等标准接口，方便各机器人之间实时数据交换，而且内置了导航、路径规划和障碍物避让等常用算法，为系统开发者节省了大量调试和集成的工作量。实际工程中，诸如TurtleBot系列和Clearpath Robotics的产品已经验证了基于ROS平台构建多机器人协同系统的可行性，这些成熟的技术方案为本系统的研发提供了坚实的实践基础。

此外，近年来以YOLO（You Only Look Once）为代表的目标检测算法，由于其端到端、实时性强的特点，正在被越来越多的实际项目采用。这种人工智能视觉技术具有轻量级，易部署等特点，可以降低部署的成本，十分适用于各种小型设备。在实际应用中，智慧校园、智能家居以及工业检测等领域已经开始引入这种的视觉识别方案，取得了明显的应用效果。%

与此同时，精准的环境感知和地图构建是多机器人协同搜索系统的重要环节。通过整合SLAM算法，如GMapping、Hector SLAM和Cartographer，以及激光雷达、RGB-D摄像头和惯性测量单元（IMU）等多种传感器的数据，系统能够实时构建并更新高精度的环境地图。这种传感器融合技术在国内外的多项机器人研究项目中已得到广泛应用，有效提升了机器人在复杂室内环境中的自主导航和定位能力。%

总的来说，结合以上的技术，可以开发基于ROS的多机器人协同搜索系统。通过整合分布式通信架构、成熟的SLAM与传感器融合技术以及YOLO等高效视觉检测算法，能够在家用、学校等小型环境中实现精准、高效的目标搜索与定位。该系统不仅为提升室内安防、智能巡检等任务的执行效率提供了有效技术支持，也为机器人群体智能和分布式控制的研究探索提供了宝贵的实践经验，预示着在未来智能家居、智慧校园及其他领域中的广阔应用前景。%

\subsection{工程基本目标}
多机器人协同搜索系统应具有两个及以上机器人，在仿真环境下使用激光雷达和光学摄像头对未知的环境进行slam建图，并且对该环境下的目标进行识别并解算坐标，完成目标在地图上的标记。与此同时，该系统应该具有自动规划路径的功能，包括避障，机器人之间避碰，未知环境探索等能力。


\section{研究思路和主要工作}
本项目将在Ubuntu20.04环境下，完成上文提及的基本目标，使用gazebo仿真技术对系统进行仿真与验证，主要工作包括以下部分：
\begin{enumerate}
\item 研究和设计一个多机器人系统架构方案，能有效地控制和协调系统中机器人之间的行为和功能逻辑关系，及设计多机器人通信架构。
\item 针对多机器人在探索过程中的目标检测问题，设计一种基于yolo模型的多机器人目标检测方案并训练调优yolo视觉识别模型。
\item 针对多机器人目标搜索定位的问题设计多传感器融合的解决方案，设计模块优化目标定位的效果。
\item 设计路径规划模块并适应多机器人协作场景，包括协同探索，碰撞规避等功能。
\end{enumerate}

\section{论文结构安排}
本论文分为以下七个章节：

第一章：绪论，对工程进行概述，介绍本工程的背景与研究思路，并且说明本论文的结构安排。

第二章：多机器人协同系统架构设计。介绍多机器人协同搜索系统的架构设计与主要模块。

第三章：机器人目标识别与检测功能设计。介绍机器人使用的目标检测模型，以及其训练调优方案。

第四章：目标检测与标记流程。说明多机器人协同搜索系统传感器融合完成目标坐标解算的流程与模块设计。

第五章：路径规划模块设计。说明多机器人协同搜索系统的路径规划模块设计与针对多机器人协作场景进行的优化。

第六章：多机器人协作探索及目标检测实验与分析。搭建仿真环境进行多机器人协作探索及目标检测，在仿真环境中进行自主探索并检测目标物，并对实验结果进行分析总结。

第七章：总结与展望。对本文已完成的研究工作和内容进行总结，分析工作中
存在的不足并对后续工作进行了展望。